% Sources verifiees WebSearch 2026-07-15 : Rth,jc TO-220 ~2-5 K/W (ROHM, fscdn.rohm.com applinote),
% Tj,max composants de puissance silicium ~150C (Renesas, STMicroelectronics datasheets IGBT),
% pont de Wheatstone quart de pont PT100 100 ohm standard (instrumexpert.com, ecotherm-mesures.fr)
% Correction 2026-08-25 (audit AUDIT_EXERCICES_2026.md) :
% - ex. 5 (ancien ex. 1) : donnees authentiques du sujet E2 2019 "Transformation d'une rame de metro"
%   (Fe1 = 14,9 kN, Fe2 = 9,2 kN EFFORTS, de = 0,86 m, rho = 0,17, eta_red = 97 %, T = 2 048 s,
%   somme = 558e6, moteurs Alsthom 96 kW / 1 035 tr/min) - questions 25 a 27 puis 28.
% - ex. 1 (ancien ex. 2) : reference corrigee "partie E" -> question D4 (E2 Disciplinaire IE 2024,
%   figure 14 et document reponse DR4).
% - ex. 3 (ancien ex. 5) : citation jury remplacee par le verbatim des elements de correction D1 du
%   Rapport_Jury_2024 ; plage -30/700 C et tolerance +/-1,5 C rattachees a la figure 13 du SUJET.
%
% Correction 2026-09-02 (complementarite fiches/exercices, protocole PROTOCOLE_COMPLEMENTARITE) :
% - ex. 1 REMPLACE : l'ancien enonce (transfert lineaire Is = f(T), -20/200 C -> 4/20 mA, application
%   T = 120 C, question D4) etait un doublon strict de la QT2 de la fiche T05
%   (fiches/T05_Phenomenes_Thermiques.tex, "QT2 - Transmetteur 4-20 mA", memes bornes, meme A.N.).
%   Le nouvel ex. 1 porte sur la SUITE de la meme partie D du meme sujet, non traitee par la fiche :
%   numerisation par le module AI 8xU/I/RTD/TC ST et mise a l'echelle PIW8 -> MW8 (questions D6 a D11).
%   Il est en outre converti en exercice docTech (encadre "Document technique -- DT10 et DT8"),
%   comblant l'absence d'exercice de lecture de document technique relevee par l'audit (§2.5).
%   Tracabilite des donnees du nouvel encadre docTech :
%     * DT10, table "Points caracteristiques maximaux" (-20 C / 200 C ; eAB mini = -1 mV /
%       eAB maxi = 10,8 mV) : VERBATIM du sujet E2_Disciplinaire_IE_2024.md, l. 942 et l. 946-949.
%     * DT10, profil du courant 4-20 mA sur -20/200 C : VERBATIM du sujet, l. 953-963 et figure 14
%       (l. 494).
%     * DT8, "2 octets par entree analogique" et "16 bits signe, signe sur le bit de poids fort" :
%       enonce du sujet l. 515 et l. 523 + elements de correction D6 du Rapport_Jury_2024.md, l. 331.
%     * DT8, ligne des courants 1,185 / 4 / 20 / 22,81 mA : VERBATIM du document reponse DR4,
%       question D7 (sujet, l. 1056-1059) ; la ligne des valeurs decimales de PIW8 y est VIDE
%       (a completer par le candidat) et le DT8 lui-meme est une image omise par la conversion.
%       La ligne PIW8 (-4 864 / 0 / 27 648 / 32 511) est donc RECONSTITUEE : la valeur pivot
%       PIW8max = 27 648 est textuellement dans le sujet (l. 539, l. 567) et dans le rapport de jury
%       (l. 373-377) ; les trois autres sont les valeurs standard de la table Siemens "Representation
%       of analog values in the current measuring ranges" pour le calibre 4-20 mA, verifiees
%       WebSearch 2026-09-02 (manuel SIMATIC S7-1500 Analog Input AI 8xU/I, p. 52 ;
%       cache.industry.siemens.com / manualslib.com).
%     * Corrige des questions 2 a 4 (a = 220/27 648 ~ 0,0080 ; b = -20 ; E = 220 ;
%       PIW8(CB=35 C) = 6 912 ; PIW8(CH=50 C) = 8 797) : VERBATIM des elements de correction
%       D9, D10 et D11 du Rapport_Jury_2024.md, l. 363-377. Consignes CB = 35 C / CH = 50 C :
%       sujet l. 545-549.
% - ex. 3 : reference du titre precisee ("epreuve ecrite disciplinaire", nom exact lu dans le sujet
%   l. 11), l'ancienne mention "Extracteur d'air" n'etant pas le nom de l'epreuve.
\expandafter\ifx\csname CAPETINCLUDEMODE\endcsname\relax
\documentclass[11pt,a4paper]{article}
% preamble.tex — packages et commandes communes a toutes les fiches CAPET SII IE
\input{glyphtounicode}
\pdfgentounicode=1
\usepackage{cmap}
\usepackage[T1]{fontenc}
\usepackage[utf8]{inputenc}
\usepackage[french]{babel}
% babel[french] active a \begin{document} des raccourcis de ponctuation (:;!?)
% qui cassent la syntaxe to[...] de circuitikz (erreur "+ or - expected") ;
% on les desactive, ce qui ne change que l'espacement fin avant ces signes.
\AtBeginDocument{\shorthandoff{:;!?}}
\usepackage[a4paper, top=1.6cm, bottom=1.6cm, left=1.4cm, right=1.4cm]{geometry}
\setlength{\headheight}{14.5pt}
\usepackage{amsmath, amssymb, amsfonts}
\usepackage{siunitx}
\sisetup{locale=FR, output-decimal-marker={,}, inter-unit-product=\ensuremath{{}\cdot{}},
         group-minimum-digits=4}
% Unites non fournies par siunitx, declarees ici une fois pour toutes (fiches + exercices).
% Les gardes \ifdefined evitent tout conflit avec une redeclaration locale restee dans un
% fichier : ne pas retirer les gardes, ne pas redeclarer ces unites dans les fichiers.
\DeclareSIUnit{\jour}{j}
\ifdefined\var\else\DeclareSIUnit{\var}{var}\fi   % puissance reactive
\ifdefined\vad\else\DeclareSIUnit{\vad}{VAD}\fi   % puissance deformante
\ifdefined\tr\else\DeclareSIUnit{\tr}{tr}\fi      % tour (tr/min)
\usepackage{array, booktabs, multirow, tabularx}
\usepackage{xcolor}
\usepackage[most, breakable]{tcolorbox}
\usepackage{multicol}
\usepackage{enumitem}
\usepackage{mathtools}
\usepackage{lmodern}
% Refonte visuelle 2026 (direction "moderne coloré") : Lato pour le texte,
% maths sans-serif assorties via newtxsf (charge apres amsmath/amssymb).
\usepackage[default]{lato}
\usepackage{newtxsf}
\usepackage{microtype}
\usepackage{fancyhdr}
\usepackage{tikz}
\usetikzlibrary{positioning,arrows.meta}
\usepackage{circuitikz}
\usepackage{lastpage}

% ── Couleurs par domaine (palette modernisée 2026) ────────────────────
\definecolor{colElec}{RGB}{29, 78, 216}        % Électrotechnique — bleu vif
\definecolor{colElecLight}{RGB}{219, 234, 254} % fond bleu clair
\definecolor{colAuto}{RGB}{21, 128, 61}        % Automatique — vert
\definecolor{colMeca}{RGB}{194, 65, 12}        % Mécanique — orange
\definecolor{colTherm}{RGB}{185, 28, 28}       % Thermique — rouge
\definecolor{colQT}{RGB}{232, 244, 233}        % (hérité, plus utilisé par questiontype)
\definecolor{colJury}{RGB}{255, 248, 225}      % fond jury (crème)
\definecolor{colJuryFrame}{RGB}{202, 138, 4}   % bordure/bandeau jury (ambre)
\definecolor{colGrayLight}{RGB}{245, 245, 247} % fond gris neutre
\definecolor{colSeparator}{RGB}{180, 180, 180} % gris séparateur
\definecolor{colExo}{RGB}{8, 145, 178}          % Banque d'exercices — cyan (distinct des fiches)
% ── Encre et gris de la refonte ───────────────────────────────────────
\definecolor{colInk}{RGB}{23, 32, 42}          % texte courant (quasi-noir)
\definecolor{colSlate}{RGB}{83, 98, 118}       % gris ardoise (méta, en-têtes)
\definecolor{colRule}{RGB}{203, 213, 225}      % filets clairs
\definecolor{colJuryBg}{RGB}{255, 248, 225}    % fond crème jury
\definecolor{colPresence}{RGB}{83, 98, 118}    % bloc 5 (présence dans les sujets)
\AtBeginDocument{\color{colInk}}

% ── Style encadré exercice + corrigé (banque d'exercices) ─────────────
\tcbset{
  exoBox/.style={
    breakable,
    enhanced,
    colback=colExo!8!white,
    colframe=colExo,
    colbacktitle=colExo,
    coltitle=white,
    fonttitle=\bfseries\small,
    left=4pt, right=4pt, top=3pt, bottom=3pt,
    before skip=5pt, after skip=5pt,
  }
}

% ── Bascule enonce / corrige (banque d'exercices) ─────────────────────
% \CAPETCORRIGE est defini en ligne de commande (pdflatex "\def\CAPETCORRIGE{1}\input{...}")
% pour compiler la variante corriges d'un meme fichier source ; par defaut (non defini),
% la compilation produit la variante enonces. \CAPETSHOW{enonce}{corrige} choisit la branche.
%
% TROIS MODES (refonte 2026-09-20, manuel de revision) :
%   - aucun drapeau   -> imprime l'enonce            (variante enonces, INCHANGE)
%   - \CAPETCORRIGE   -> imprime le corrige          (variante corriges, INCHANGE)
%   - \CAPETMANUEL    -> imprime l'enonce ET met le corrige de cote pour l'annexe
%
% Le mode manuel stocke chaque corrige dans une macro (\csgdef), sans ecriture
% fichier ni re-tokenisation : le contenu (circuitikz, accents, %) traverse intact.
% L'annexe les rejoue ensuite dans l'ordre via \CAPETANNEXE.
\usepackage{etoolbox}

\newcounter{CAPETexo}          % numero continu d'exercice dans le manuel
\setcounter{CAPETexo}{0}
\newcounter{CAPETannexe}       % compteur de relecture en annexe

% Titre courant de la notion en cours (pose par \CAPETNOTION), sert d'intertitre annexe
\newcommand{\CAPETnotioncode}{}
\newcommand{\CAPETnotiontitre}{}

% \CAPETNOTION{code}{titre} : ouvre une notion dans le manuel (memorise le libelle,
% et marque le rang du premier exercice de la notion pour l'intertitre d'annexe).
\newcommand{\CAPETNOTION}[2]{%
  \renewcommand{\CAPETnotioncode}{#1}%
  \renewcommand{\CAPETnotiontitre}{#2}%
  \csgdef{CAPETnotionstart\the\numexpr\value{CAPETexo}+1\relax}{#1 -- #2}%
}

% ── \degree robuste ───────────────────────────────────────────────────
% Certains fichiers d'exercices utilisent \degree. Ni le preambule ni siunitx
% ne le definissent, et \providecommand local ne suffit pas : sous hyperref,
% \degree existe deja comme commande factice (\ERROR), donc \providecommand
% s'abstient et l'appel echoue a la compilation du manuel (erreur invisible
% en compilation isolee). On le (re)definit donc inconditionnellement ici,
% apres le chargement des packages, pour tout le corpus.
\AtBeginDocument{\providecommand{\degree}{}\renewcommand{\degree}{\ensuremath{^\circ}}}

% Libelle decoratif : les fichiers d'exercices utilisent aussi \CAPETSHOW dans
% leur \fancyhead et leur bloc de titre, pour afficher "(enonces)" ou "(corriges)".
% Ces appels-la ne sont PAS des exercices : en mode manuel ils ne doivent ni
% consommer un numero, ni alimenter l'annexe. On les reconnait a leur premier
% argument, strictement egal au libelle ci-dessous.
\newcommand{\CAPETLIBELLEENONCES}{(\'enonc\'es)}

\ifdefined\CAPETMANUEL
  % --- Mode manuel : enonce en place, corrige differe vers l'annexe ---
  \newcommand{\CAPETSHOW}[2]{%
    \def\CAPETargun{#1}%
    \ifx\CAPETargun\CAPETLIBELLEENONCES
      #1%            appel decoratif : rendu tel quel, hors numerotation
    \else
      \CAPETSHOWEXO{#1}{#2}%
    \fi
  }
  \newcommand{\CAPETSHOWEXO}[2]{%
    \stepcounter{CAPETexo}%
    \csgdef{CAPETcorrige\theCAPETexo}{#2}%
    \edef\CAPETtmpnum{\theCAPETexo}%
    #1%
    \par\vspace{2pt}%
    {\footnotesize\itshape\color{colExo}%
     \hyperref[capet:cor:\CAPETtmpnum]{Corrige de l'exercice \CAPETtmpnum\ $\rightarrow$ p.~\pageref{capet:cor:\CAPETtmpnum}}%
     \label{capet:eno:\CAPETtmpnum}}%
  }
  % \CAPETANNEXE : rejoue tous les corriges captures, dans l'ordre, avec intertitres
  \newcommand{\CAPETANNEXE}{%
    \setcounter{CAPETannexe}{0}%
    \loop\ifnum\value{CAPETannexe}<\value{CAPETexo}%
      \stepcounter{CAPETannexe}%
      \edef\CAPETcurnum{\theCAPETannexe}%
      \ifcsname CAPETnotionstart\CAPETcurnum\endcsname
        \par\medskip
        {\normalsize\bfseries\color{colExo}\csuse{CAPETnotionstart\CAPETcurnum}}%
        \par\nobreak\smallskip
      \fi
      \begin{tcolorbox}[exoBox, breakable,
        title={Corrige de l'exercice \CAPETcurnum\
               \hfill\normalfont\footnotesize
               \hyperref[capet:eno:\CAPETcurnum]{enonce p.~\pageref{capet:eno:\CAPETcurnum}}}]%
      \label{capet:cor:\CAPETcurnum}%
      \csuse{CAPETcorrige\CAPETcurnum}%
      \end{tcolorbox}%
    \repeat
  }
\else
  % --- Modes historiques, comportement strictement inchange ---
  \newcommand{\CAPETSHOW}[2]{\ifdefined\CAPETCORRIGE #2\else #1\fi}
  \newcommand{\CAPETANNEXE}{}
\fi

% ── Difficulté 1 à 3 (étoiles pleines, amssymb) ───────────────────────
\newcommand{\difficulte}[1]{%
  \ifcase#1\relax
  \or $\bigstar$\;niveau 1/3%
  \or $\bigstar\bigstar$\;niveau 2/3%
  \or $\bigstar\bigstar\bigstar$\;niveau 3/3%
  \fi
}

% ── Style encadré extrait de document technique (banque d'exercices) ──
\tcbset{
  docTech/.style={
    breakable,
    enhanced,
    colback=colGrayLight,
    colframe=colSeparator,
    fonttitle=\bfseries\small\itshape,
    coltitle=black,
    colbacktitle=colGrayLight,
    left=4pt, right=4pt, top=3pt, bottom=3pt,
    before skip=4pt, after skip=4pt,
  }
}

% ── Style tcolorbox de section principale ────────────────────────────
% Refonte 2026 : bandeau titre coloré arrondi + fond très pâle, sans cadre.
\tcbset{
  ficheSection/.style={
    breakable, enhanced jigsaw,
    frame hidden, boxrule=0pt, arc=1.8mm,
    colback=#1!4!white,
    colbacktitle=#1,
    coltitle=white,
    fonttitle=\sffamily\bfseries,
    left=7pt, right=7pt, top=5pt, bottom=5pt,
    toptitle=3pt, bottomtitle=3pt,
    before skip=10pt, after skip=8pt,
  }
}

% ── Style encadré question type ──────────────────────────────────────
% Carte blanche, bandeau sombre, liseré très clair.
\tcbset{
  questiontype/.style={
    breakable, enhanced jigsaw,
    colback=white,
    frame hidden, boxrule=0pt, arc=1.8mm,
    colbacktitle=colElec!78!black,
    coltitle=white,
    fonttitle=\sffamily\bfseries\small,
    left=7pt, right=7pt, top=4pt, bottom=4pt,
    toptitle=2.5pt, bottomtitle=2.5pt,
    before skip=8pt, after skip=8pt,
    borderline={0.6pt}{0pt}{colElec!12!white},
  }
}

% ── Style encadré jury ───────────────────────────────────────────────
% Bandeau ambre, fond crème, sans cadre.
\tcbset{
  jury/.style={
    breakable, enhanced jigsaw,
    colback=colJuryBg,
    frame hidden, boxrule=0pt, arc=1.8mm,
    colbacktitle=colJuryFrame,
    coltitle=white,
    fonttitle=\sffamily\bfseries,
    left=7pt, right=7pt, top=4pt, bottom=4pt,
    toptitle=3pt, bottomtitle=3pt,
    before skip=10pt, after skip=8pt,
  }
}

% ── En-tête / pied de page communs (refonte 2026) ─────────────────────
% #1 code fiche, #2 titre court, #3 couleur domaine
\newcommand{\ficheHeaders}[3]{%
  \fancyhead[L]{\small\sffamily{\bfseries\textcolor{#3}{#1}}\hspace{0.7em}\textcolor{colSlate}{#2}}%
  \fancyhead[R]{\small\sffamily\textcolor{colSlate}{CAPET SII IE\, --- \,Session 2026}}%
  \fancyfoot[C]{\small\sffamily\textcolor{colSlate}{\thepage\,/\,\pageref{LastPage}}}%
  \fancyhead[C]{}\fancyfoot[L]{}\fancyfoot[R]{}%
}

% ── Bloc en-tête de fiche (BLOC 0, refonte 2026) ──────────────────────
% #1 code  #2 titre  #3 couleur domaine  #4 domaine  #5 ligne présence  #6 sous-thèmes
\newcommand{\ficheEntete}[6]{%
  \begin{tcolorbox}[enhanced, breakable,
    frame hidden, boxrule=0pt, arc=2.5mm,
    interior style={left color=#3, right color=#3!75!black},
    left=9pt, right=9pt, top=7pt, bottom=7pt,
    before skip=0pt, after skip=6pt]
    {\color{white}\sffamily\bfseries\LARGE #1\quad #2}\\[4pt]
    {\color{#3!22!white}\sffamily\small \MakeUppercase{#4}\quad\textbar\quad #5}\\[1pt]
    {\color{#3!35!white}\sffamily\footnotesize Sous-thèmes : #6}
  \end{tcolorbox}%
}

% ── Mise en page ─────────────────────────────────────────────────────
\pagestyle{fancy}
\fancyhf{}
\renewcommand{\headrulewidth}{0.5pt}
\renewcommand{\footrulewidth}{0pt}
\AtBeginDocument{%
  \renewcommand{\headrule}{{\color{colRule}\hrule width\headwidth height 0.5pt}}%
}

\setlist[itemize]{leftmargin=1.2em, itemsep=1pt, topsep=2pt, parsep=0pt}
\setlist[enumerate]{leftmargin=1.2em, itemsep=3pt, topsep=2pt, parsep=0pt}

\setlength{\parskip}{2pt}
\setlength{\columnsep}{10pt}

\def\CAPETBEGINFILE{\begin{document}}
\def\CAPETENDFILE{\end{document}}
\else
\def\CAPETBEGINFILE{}
\def\CAPETENDFILE{}
\fi

\ifdefined\tr\else\DeclareSIUnit{\tr}{tr}\fi

\fancyhead[L]{\textbf{\textcolor{colExo}{T05}} - Exercices - Ph\'enom\`enes thermiques}
\fancyhead[R]{\small\textit{Banque d'exercices \CAPETSHOW{(\'enonc\'es)}{(corrig\'es)}}}
\fancyfoot[C]{\small\thepage\,/\,\pageref{LastPage}}

\CAPETBEGINFILE

\begin{tcolorbox}[
  enhanced, breakable,
  colback=colExo,
  colframe=colExo,
  left=6pt, right=6pt, top=5pt, bottom=5pt,
  before skip=0pt, after skip=4pt
]
  {\color{white}\bfseries\large T05 \quad Ph\'enom\`enes thermiques}\\[2pt]
  {\color{white}\small Banque d'exercices -- 5 exercices \CAPETSHOW{(\'enonc\'es)}{(corrig\'es)}}
\end{tcolorbox}

\begin{tcolorbox}[exoBox,
  title={Exercice 1 -- \difficulte{1} \hfill\normalfont\footnotesize 2024, \'epreuve \'ecrite disciplinaire, questions D6 \`a D11}
]
\CAPETSHOW{
Dans l'armoire \'electrique d'un extracteur d'air, la temp\'erature $T_{ARM}$ est mesur\'ee par un thermocouple de type J, mis en forme par un transmetteur \guillemotleft\ TT~\guillemotright\ (r\'ef\'erence KOS~839) dont la sortie attaque le module d'entr\'ees analogiques \emph{AI 8xU/I/RTD/TC ST} d'un automate S7-1500. Le module recopie le courant re\c cu dans le mot d'entr\'ee $PIW8$ ; un programme de \guillemotleft\ mise \`a l'\'echelle~\guillemotright\ en d\'eduit le mot interne $MW8$, image de la temp\'erature en \si{\celsius}.

\begin{tcolorbox}[docTech, title={Document technique -- DT10 (profils capteur / transmetteur) et DT8 (module d'entr\'ees analogiques)}]
\textbf{DT10 -- Profil de la tension fournie par le capteur (thermocouple type J).}

\begin{center}
\renewcommand{\arraystretch}{1.15}
\begin{tabular}{lcc}
\toprule
\textbf{Points caract\'eristiques maximaux} & & \\
\midrule
Temp\'erature mesur\'ee (\si{\celsius}) & $-20$ & $200$ \\
Tension d\'elivr\'ee par le capteur (\si{\milli\volt}) & $e_{AB\,mini} = -1$ & $e_{AB\,maxi} = 10{,}8$ \\
\bottomrule
\end{tabular}
\end{center}

\textbf{DT10 -- Profil du courant d\'elivr\'e par le transmetteur TT} : $I_s$ varie lin\'eairement de $I_{s\,mini} = \SI{4}{\milli\ampere}$ \`a $I_{s\,maxi} = \SI{20}{\milli\ampere}$ lorsque $T_{ARM}$ varie de $-20$ \`a $\SI{200}{\celsius}$.

\smallskip
\textbf{DT8 -- R\'esolution des valeurs analogiques ($PIW_i$).} Une entr\'ee analogique occupe une adresse m\'emoire de 2 octets. Repr\'esentation des valeurs analogiques dans les calibres de mesure en courant :

\begin{center}
\renewcommand{\arraystretch}{1.15}
\begin{tabular}{lcccc}
\toprule
Courant mesur\'e (\si{\milli\ampere}) & $1{,}185$ & $4$ & $20$ & $22{,}81$ \\
Valeur d\'ecimale du mot $PIW8$ & $-4\,864$ & $0$ & $27\,648$ & $32\,511$ \\
\bottomrule
\end{tabular}
\end{center}
\end{tcolorbox}

\begin{enumerate}[leftmargin=1.5em]
  \item \textbf{Relever} dans le document technique DT8 le nombre d'octets sur lequel est cod\'ee une entr\'ee analogique, en d\'eduire le nombre de bits, puis \textbf{pr\'eciser} le r\^ole du bit de poids fort. \textbf{Relever} \'egalement les valeurs d\'ecimales de $PIW8$ associ\'ees \`a $I_s = \SI{4}{\milli\ampere}$ et $I_s = \SI{20}{\milli\ampere}$.
  \item Le mot $MW8$ doit contenir la temp\'erature en \si{\celsius}. \'Etablir l'\'equation de la droite $MW8 = f(PIW8)$ sous la forme $MW8 = a\cdot PIW8 + b$, en exprimant $a$ et $b$ \`a l'aide de $MW8_{min}$, $MW8_{max}$ et $PIW8_{max}$. Faire l'application num\'erique.
  \item On note $E$ l'\'etendue de mesure du transmetteur. Donner sa valeur et \textbf{montrer} que l'\'equation de mise \`a l'\'echelle s'\'ecrit
  \[ MW8 = \dfrac{E}{27\,648}\cdot PIW8 - 20 \]
  \item La ventilation de l'armoire est command\'ee par deux consignes de temp\'erature : $C_B = \SI{35}{\celsius}$ (petite vitesse) et $C_H = \SI{50}{\celsius}$ (grande vitesse). Ces seuils doivent \^etre compar\'es \emph{dans l'automate} \`a la valeur brute $PIW8$. Calculer les deux valeurs de $PIW8$ correspondantes.
  \item En exploitant la premi\`ere ligne du tableau DT8, indiquer ce que traduit une valeur $PIW8$ n\'egative et l'usage que le programme peut en faire.
\end{enumerate}

}{
\textbf{Corrig\'e.}

\begin{enumerate}[leftmargin=1.5em]
  \item \textbf{Extraction.} DT8 indique 2 octets par entr\'ee analogique, soit $2\times8 = 16$ bits. Le codage \'etant sign\'e, le bit de poids fort porte le signe : la r\'esolution utile n'est donc que de 15 bits. Le tableau DT8 donne directement $PIW8_{min} = 0$ pour $I_s = \SI{4}{\milli\ampere}$ et $PIW8_{max} = 27\,648$ pour $I_s = \SI{20}{\milli\ampere}$.
  \item Aux extr\'emit\'es de la plage : $(PIW8 = 0 \Rightarrow MW8_{min} = -20)$ et $(PIW8 = 27\,648 \Rightarrow MW8_{max} = 200)$. L'ordonn\'ee \`a l'origine se lit imm\'ediatement : $b = MW8_{min} = -20$. La pente s'obtient avec le second point :
  \[ a = \dfrac{MW8_{max} - MW8_{min}}{PIW8_{max}} = \dfrac{200 - (-20)}{27\,648} = \dfrac{220}{27\,648} \approx 0{,}0080 \]
  \[ MW8 = \dfrac{220}{27\,648}\cdot PIW8 - 20 \]
  \item L'\'etendue de mesure est l'\'ecart entre les temp\'eratures extr\^emes du transmetteur :
  \[ E = T_{max} - T_{min} = 200 - (-20) = \SI{220}{\celsius} \]
  En rempla\c cant le num\'erateur $220$ par $E$ dans le r\'esultat de la question pr\'ec\'edente, on obtient bien la forme demand\'ee $MW8 = \dfrac{E}{27\,648}\cdot PIW8 - 20$. \'Ecrite ainsi, l'\'equation reste valable si l'on rechange l'\'etalonnage du transmetteur : seule la constante $E$ est \`a mettre \`a jour dans le programme.
  \item On isole $PIW8$ :
  \[ MW8 = \dfrac{E}{27\,648}\cdot PIW8 - 20 \iff PIW8 = \dfrac{27\,648}{E}\,(MW8 + 20) \]
  \[ PIW8(C_B) = \dfrac{27\,648}{220}\,(35+20) = 125{,}67\times55 = 6\,912 \]
  \[ PIW8(C_H) = \dfrac{27\,648}{220}\,(50+20) = 125{,}67\times70 \approx 8\,797 \]
  \item Le tableau DT8 associe $I_s = \SI{1,185}{\milli\ampere}$ \`a $PIW8 = -4\,864$ : le module code en n\'egatif les courants inf\'erieurs \`a \SI{4}{\milli\ampere}, c'est-\`a-dire hors de la plage normalis\'ee. Un $PIW8$ n\'egatif ne correspond donc pas \`a une temp\'erature inf\'erieure \`a $\SI{-20}{\celsius}$, mais \`a une rupture de la boucle de courant (capteur d\'ebranch\'e, fil coup\'e, transmetteur hors service). Le programme peut tester $PIW8 < 0$ (et sym\'etriquement $PIW8 > 27\,648$) pour lever un bit \texttt{defaut\_mesure} avant tout calcul de mise \`a l'\'echelle.
\end{enumerate}

\smallskip
\textbf{Conclusion.} Comparer directement $PIW8$ aux seuils $6\,912$ et $8\,797$ \'evite \`a l'automate toute division en ligne : la commande du ventilateur se r\'eduit \`a deux comparaisons d'entiers, plus rapides et exemptes d'erreur d'arrondi que le calcul de $MW8$. Surtout, l'exigence de s\^uret\'e est port\'ee par le codage lui-m\^eme : la plage $[-4\,864\,;\,32\,511]$ plus large que $[0\,;\,27\,648]$ est ce qui rend le d\'efaut de mesure \emph{d\'etectable}, l\`a o\`u une plage born\'ee \`a $[0\,;\,27\,648]$ confondrait une rupture de boucle avec une temp\'erature basse.
}
\end{tcolorbox}

\begin{tcolorbox}[exoBox,
  title={Exercice 2 -- \difficulte{2} \hfill\normalfont\footnotesize exercice in\'edit}
]
\CAPETSHOW{
Un transistor de puissance en bo\^itier TO-220 dissipe une puissance $P = \SI{8}{\watt}$ en r\'egime permanent. Il est mont\'e sur un dissipateur (radiateur) refroidi par convection naturelle. Les r\'esistances thermiques du montage sont : jonction-bo\^itier $R_{th,jc} = \SI{2}{\kelvin\per\watt}$, bo\^itier-dissipateur (interface, pastille isolante) $R_{th,cs} = \SI{0,5}{\kelvin\per\watt}$, dissipateur-ambiant $R_{th,sa} = \SI{5}{\kelvin\per\watt}$. La temp\'erature ambiante est $T_{amb} = \SI{35}{\celsius}$ et le fabricant sp\'ecifie une temp\'erature de jonction maximale $T_{j,max} = \SI{150}{\celsius}$.

\begin{center}
\begin{circuitikz}[european, scale=0.85]
\draw
(0,0) to[R, l=$R_{th,jc}$] (2.5,0)
to[R, l=$R_{th,cs}$] (5,0)
to[R, l=$R_{th,sa}$] (7.5,0)
;
\draw (0,-1.3) to[I, l=$P$] (0,0);
\draw[-latex] (0,-1.8) -- (7.5,-1.8);
\node at (0,-2.1) {$T_j$};
\node at (2.5,-2.1) {$T_c$};
\node at (5,-2.1) {$T_s$};
\node at (7.5,-2.1) {$T_{amb}$};
\draw[dotted] (0,0) -- (0,-1.8);
\draw[dotted] (2.5,0) -- (2.5,-1.8);
\draw[dotted] (5,0) -- (5,-1.8);
\draw[dotted] (7.5,0) -- (7.5,-1.8);
\end{circuitikz}
\end{center}

\begin{enumerate}[leftmargin=1.5em]
  \item Justifier l'analogie entre ce sch\'ema thermique (source de puissance $P$, r\'esistances thermiques en s\'erie, temp\'eratures aux n\oe uds) et un circuit \'electrique.
  \item Calculer la r\'esistance thermique totale jonction-ambiant $R_{th,ja}$, puis la temp\'erature de jonction $T_j$ en r\'egime permanent.
  \item Conclure quant au respect de $T_{j,max}$, et proposer une action si la marge est insuffisante.
\end{enumerate}

}{
\textbf{Corrig\'e.}

\begin{enumerate}[leftmargin=1.5em]
  \item L'analogie \'electrique-thermique associe la puissance dissip\'ee $P$ (source) au courant $I$, la temp\'erature $T$ au potentiel $V$, et la r\'esistance thermique $R_{th}$ (en $\si{\kelvin\per\watt}$) \`a la r\'esistance \'electrique (en $\si{\ohm}$). La loi $\Delta T = R_{th}\cdot P$ est l'analogue direct de la loi d'Ohm $\Delta V = R\cdot I$ ; des r\'esistances thermiques travers\'ees par le m\^eme flux se comportent comme des r\'esistances \'electriques en s\'erie.
  \item Les trois r\'esistances thermiques sont travers\'ees par le m\^eme flux $P$ (m\^eme chemin d'\'evacuation de la chaleur), donc en s\'erie :
  \[ R_{th,ja} = R_{th,jc} + R_{th,cs} + R_{th,sa} = 2 + 0{,}5 + 5 = \SI{7,5}{\kelvin\per\watt} \]
  \[ T_j = T_{amb} + R_{th,ja} \cdot P = 35 + 7{,}5 \times 8 = 35 + 60 = \SI{95}{\celsius} \]
  \item $T_j = \SI{95}{\celsius} < T_{j,max} = \SI{150}{\celsius}$ : la marge est de $\SI{55}{\celsius}$, le point de fonctionnement est s\^ur. Si la marge \'etait insuffisante, on r\'eduirait $R_{th,sa}$ (dissipateur plus grand ou ventilation forc\'ee) ou $R_{th,cs}$ (meilleure interface thermique), ces deux termes \'etant les plus facilement modifiables sans changer le composant.
\end{enumerate}

\smallskip
\textbf{Conclusion.} Le crit\`ere de dimensionnement thermique $T_j \leq T_{j,max}$ est v\'erifi\'e avec $\SI{55}{\celsius}$ de marge ; comme $R_{th,sa}$ repr\'esente \`a elle seule $5/7{,}5 = 67\,\%$ de la r\'esistance totale, c'est le choix du dissipateur qui pilote la temp\'erature de jonction et donc la dur\'ee de vie du composant.
}
\end{tcolorbox}

\begin{tcolorbox}[exoBox,
  title={Exercice 3 -- \difficulte{2} \hfill\normalfont\footnotesize d'apr\`es 2024, \'epreuve \'ecrite disciplinaire, question D1}
]
\CAPETSHOW{
Un thermocouple de type J (couple Fer/Constantan) surveille la temp\'erature d'un moteur \'electrique. Sa jonction chaude (jonction de mesure) est plong\'ee dans le bobinage \`a temp\'erature $T_{mes}$, et sa jonction froide (jonction de r\'ef\'erence, au niveau du conditionneur) est \`a temp\'erature $T_{ref}$, non n\'ecessairement nulle. Le coefficient de Seebeck moyen du couple, relev\'e dans la documentation constructeur (il n'est pas fourni par le sujet), vaut $k_S \approx \SI{51}{\micro\volt\per\kelvin}$.

\begin{center}
\begin{circuitikz}[european, scale=0.85]
\draw (0,3) to[short, *-] (2,3);
\draw (0,0) to[short, *-] (2,0);
\draw (2,3) to[cute inductor, l=fer] (5,1.5) to[cute inductor, l=constantan] (2,0);
\node[circ] at (5,1.5) {};
\node[right] at (5,1.5) {jonction chaude, $T_{mes}$};
\draw (0,3) to[voltmeter] (0,0);
\node[left] at (-0.5,1.5) {jonction froide, $T_{ref}$};
\end{circuitikz}
\end{center}

\begin{enumerate}[leftmargin=1.5em]
  \item Nommer le ph\'enom\`ene physique exploit\'e par le thermocouple, et donner l'expression de la f.\'e.m. $e_{AB}$ g\'en\'er\'ee.
  \item La jonction froide n'est pas maintenue \`a $\SI{0}{\celsius}$ mais \`a $T_{ref} = \SI{22}{\celsius}$ (temp\'erature ambiante du bo\^itier). Expliquer pourquoi une mesure directe de $e_{AB}$ sans compensation fausserait la lecture de $T_{mes}$, et proposer une solution.
  \item Calculer la f.\'e.m. mesur\'ee pour $T_{mes} = \SI{180}{\celsius}$ et $T_{ref} = \SI{22}{\celsius}$.
\end{enumerate}

}{
\textbf{Corrig\'e.}

\begin{enumerate}[leftmargin=1.5em]
  \item C'est l'effet Seebeck : la jonction de deux m\'etaux diff\'erents ($A$ = fer, $B$ = constantan) port\'ee \`a deux temp\'eratures distinctes g\'en\`ere une f.\'e.m. proportionnelle \`a l'\'ecart de temp\'erature :
  \[ e_{AB} = k_S \cdot (T_{mes} - T_{ref}) \]
  \item Si l'on suppose \`a tort $T_{ref} = \SI{0}{\celsius}$ alors que la jonction froide est en r\'ealit\'e \`a $\SI{22}{\celsius}$, la f.\'e.m. mesur\'ee correspond \`a un \'ecart $(T_{mes}-22)$ et non $(T_{mes}-0)$ : la temp\'erature d\'eduite serait sous-estim\'ee de $\SI{22}{\celsius}$. La solution classique est la compensation de soudure froide : un capteur local (souvent \`a semi-conducteur) mesure $T_{ref}$ et le conditionneur ajoute \'electroniquement le terme correctif $k_S \cdot T_{ref}$.
  \item $e_{AB} = 51\times10^{-6}\times(180-22) = 51\times10^{-6}\times158 \approx \SI{8,06}{\milli\volt}$.
\end{enumerate}

\smallskip
\textbf{Conclusion.} Pour un \'ecart de $\SI{158}{\celsius}$, la f.\'e.m. utile ne vaut que $\SI{8,06}{\milli\volt}$ : elle est bien \guillemotleft\ de l'ordre du millivolt~\guillemotright, donc inexploitable directement par l'entr\'ee analogique de l'automate. Elle impose la cha\^ine de conditionnement d\'ecrite par le sujet -- compensation de soudure froide, amplification, puis conversion tension $\rightarrow$ courant $4$--$\SI{20}{\milli\ampere}$ assur\'ee par le transmetteur KOS 839.

\smallskip
\textbf{Conseil de m\'ethode.} Le choix d'un type de thermocouple se justifie en confrontant la grandeur \`a mesurer \`a la plage conseill\'ee et \`a la tol\'erance publi\'ees par le constructeur : la figure 13 du sujet 2024 annonce, pour le type J (Fer/Constantan), une plage de $-30$ \`a $\SI{700}{\celsius}$ et une tol\'erance de $\pm\SI{1,5}{\celsius}$, puis fait arbitrer sur le prix et le milieu d'utilisation.

\smallskip
\textit{Point de vigilance jury (rapport de jury CAPET 2024, \'el\'ements de correction de la question D1) : \guillemotleft\ Le type T est le capteur avec une meilleure tol\'erance et une r\'ep\'etitivit\'e remarquable mais c'est aussi le plus cher. Utilis\'e dans un milieu non humide, comme notre cas, le capteur J convient car moins cher \guillemotright\ -- la justification attendue croise donc trois crit\`eres (plage, tol\'erance, co\^ut et milieu), et non le seul crit\`ere de plage.}
}
\end{tcolorbox}

\begin{tcolorbox}[exoBox,
  title={Exercice 4 -- \difficulte{2} \hfill\normalfont\footnotesize exercice in\'edit}
]
\CAPETSHOW{
Une sonde PT100 est ins\'er\'ee dans le bras d'un pont de Wheatstone en montage quart de pont : trois r\'esistances fixes $R = \SI{100}{\ohm}$ et la PT100 $R_T$, aliment\'ees par une tension continue $E = \SI{5}{\volt}$. Le point $B$ est pris comme r\'ef\'erence de potentiel ($V_B = 0$) et le point $A$ est port\'e au potentiel $E$. La caract\'eristique lin\'earis\'ee de la sonde est $R_T = R_0(1+\alpha T)$ avec $R_0 = \SI{100}{\ohm}$ et $\alpha = 3{,}85\times10^{-3}\ \si{\per\kelvin}$.

\begin{center}
\begin{circuitikz}[european, scale=0.85]
\draw
(0,4) node[above]{$A$}
to[R, l=$R$] (-3,2) node[left]{$V_+$}
to[R, l=$R$] (0,0) node[above right]{$B$}
;
\draw
(0,4) to[R, l=$R$] (3,2) node[right]{$V_-$}
to[R, l=$R_T$] (0,0)
;
\draw (-3,2) to[voltmeter] (3,2);
\draw (0,4) -- (-5,4) to[V, l=$E$] (-5,0) -- (0,0);
\draw (0,0) -- (0,-0.7) node[ground]{};
\node at (0,-2.3) {$U_{pont} = V_+ - V_-$};
\end{circuitikz}
\end{center}

\begin{enumerate}[leftmargin=1.5em]
  \item \`A $T = \SI{0}{\celsius}$, montrer que le pont est \'equilibr\'e ($U_{pont} = 0$).
  \item \'Etablir l'expression \emph{exacte} de $U_{pont}$ en fonction de $E$, $R$ et $R_T$, puis la lin\'eariser dans l'approximation $\Delta R = R_T - R \ll R$. Pr\'eciser le signe obtenu et l'interpr\'eter.
  \item Calculer $U_{pont}$ pour $T = \SI{80}{\celsius}$ avec les deux expressions (exacte puis lin\'earis\'ee), et commenter l'\'ecart.
\end{enumerate}

}{
\textbf{Corrig\'e.}

\begin{enumerate}[leftmargin=1.5em]
  \item \`A $T=\SI{0}{\celsius}$, $R_T = R_0(1+0) = R_0 = \SI{100}{\ohm} = R$ : les quatre branches du pont sont identiques deux \`a deux, le pont est \'equilibr\'e et $U_{pont}=0$. C'est la configuration de r\'ef\'erence du montage.
  \item Les deux branches sont deux diviseurs de tension aliment\'es sous $E$ :
  \[ V_+ = E\cdot\dfrac{R}{R+R} = \dfrac{E}{2}, \qquad V_- = E\cdot\dfrac{R_T}{R+R_T} \]
  D'o\`u, apr\`es mise au m\^eme d\'enominateur, l'expression exacte :
  \[ U_{pont} = V_+ - V_- = E\cdot\dfrac{(R+R_T) - 2R_T}{2(R+R_T)} = E\cdot\dfrac{R - R_T}{2\,(R+R_T)} = -\,E\cdot\dfrac{\Delta R}{2\,(2R+\Delta R)} \]
  Pour $\Delta R \ll R$, on n\'eglige $\Delta R$ devant $2R$ au d\'enominateur :
  \[ U_{pont} \approx -\,E\cdot\dfrac{\Delta R}{4R} = -\,E\cdot\dfrac{R_0\,\alpha\,T}{4R} \quad\text{soit, avec } R = R_0,\quad U_{pont} \approx -\,\dfrac{E\,\alpha\,T}{4} \]
  Le signe est \textbf{n\'egatif} pour $T > \SI{0}{\celsius}$ : la r\'esistance de la sonde augmente avec la temp\'erature, donc $V_-$ (diviseur $R/R_T$) d\'epasse $V_+ = E/2$, qui reste fixe. La sensibilit\'e du pont vaut $\mathrm{d}U_{pont}/\mathrm{d}T = -E\alpha/4 = -\SI{4,81}{\milli\volt\per\kelvin}$.
  \item \`A $T = \SI{80}{\celsius}$ : $\Delta R = R_0\,\alpha\,T = 100\times3{,}85\times10^{-3}\times80 = \SI{30,8}{\ohm}$, soit $R_T = \SI{130,8}{\ohm}$.

  \emph{Expression exacte} : $V_+ = \SI{2,5}{\volt}$ et $V_- = 5\times\dfrac{130{,}8}{230{,}8} = \SI{2,834}{\volt}$, donc
  \[ U_{pont} = 2{,}500 - 2{,}834 = \SI{-0,334}{\volt} \]

  \emph{Expression lin\'earis\'ee} : $U_{pont} \approx -\dfrac{5 \times 3{,}85\times10^{-3}\times 80}{4} = -\dfrac{1{,}54}{4} = \SI{-0,385}{\volt}$.

  L'\'ecart relatif atteint $(0{,}385-0{,}334)/0{,}334 \approx 15\,\%$ : \`a $\SI{80}{\celsius}$, $\alpha T = 0{,}308$, donc $\Delta R/R \approx 31\,\%$ et l'hypoth\`ese $\Delta R \ll R$ n'est plus l\'egitime.
\end{enumerate}

\smallskip
\textbf{Conclusion.} Avec une sensibilit\'e de $\SI{4,81}{\milli\volt\per\kelvin}$, le pont d\'elivre au plus quelques centaines de millivolts sur l'\'etendue de mesure : le signal doit \^etre amplifi\'e (amplificateur d'instrumentation) avant num\'erisation. Surtout, l'erreur de lin\'earisation de $15\,\%$ \`a $\SI{80}{\celsius}$ est sans commune mesure avec la classe de pr\'ecision vis\'ee par une cha\^ine de mesure industrielle (typiquement $0{,}5\,\%$) : au-del\`a de quelques dizaines de degr\'es, c'est l'expression exacte -- ou une lin\'earisation par morceaux dans le conditionneur -- qu'il faut retenir.
}
\end{tcolorbox}

\begin{tcolorbox}[exoBox,
  title={Exercice 5 -- \difficulte{3} \hfill\normalfont\footnotesize 2019, sujet E2 -- rame de m\'etro, questions 25 \`a 28}
]
\CAPETSHOW{
Une rame de m\'etro de maintenance est tract\'ee par des moteurs \`a courant continu. Chaque moteur entra\^ine un essieu par l'interm\'ediaire d'un r\'educteur de rapport de r\'eduction $\rho = 0{,}17$ et de rendement $\eta_{red} = 97\,\%$ ; deux roues de diam\`etre $d_e = \SI{0,86}{\meter}$ sont solidaires de l'essieu.

L'effort de traction $F_e$ que doit produire un moteur n'est pas constant sur le trajet : le trajet aller-retour, phase sur chantier comprise, se d\'ecompose en 18 phases \`a effort constant, pour une dur\'ee totale $T = 2\,048\ \si{\second}$. On donne $F_{e1} = \SI{14,9}{\kilo\newton}$ pendant $\Delta t_1 = \SI{42}{\second}$, puis $F_{e2} = \SI{9,2}{\kilo\newton}$ pendant $\Delta t_2 = \SI{51}{\second}$ ; les seize phases restantes contribuent pour un terme cumul\'e $\displaystyle\sum_{i=3}^{18} C_{mi}^2 \cdot \Delta t_i = 558\times10^{6}\ \si{\newton\squared\meter\squared\second}$.

Le couple thermique \'equivalent $C_{th}$ est le couple constant qui provoquerait le m\^eme \'echauffement du moteur sur le cycle :
\[ C_{th} = \sqrt{\dfrac{1}{T}\left( C_{m1}^2\,\Delta t_1 + C_{m2}^2\,\Delta t_2 + \sum_{i=3}^{18} C_{mi}^2\,\Delta t_i \right)} \]

\begin{center}
\begin{tikzpicture}[scale=0.85, >=Stealth]
  \draw[-latex] (0,0) -- (8.7,0) node[right]{$t$};
  \draw[-latex] (0,0) -- (0,2.6) node[above]{$F_e(t)$};
  \draw[dashed, gray] (0,2) -- (2.4,2); \node[left] at (0,2) {$F_{e1}$};
  \draw[dashed, gray] (0,1.2) -- (5.5,1.2); \node[left] at (0,1.2) {$F_{e2}$};
  \draw[thick]
    (0,2) -- (2.4,2) -- (2.4,1.2) -- (5.5,1.2) -- (5.5,0.5);
  \draw[thick, dotted] (5.5,0.5) -- (8,0.5);
  \draw[dotted] (2.4,0) -- (2.4,2);
  \node[below] at (2.4,0) {$\Delta t_1$};
  \draw[dotted] (5.5,0) -- (5.5,1.2);
  \node[below] at (5.5,0) {$\Delta t_1{+}\Delta t_2$};
  \node at (7,1.9) {\footnotesize phases 3 \`a 18};
  \node[below] at (0,0) {$0$};
  \draw[dotted] (8,0) -- (8,0.5);
  \node[below] at (8,0) {$T$};
\end{tikzpicture}
\end{center}

\begin{enumerate}[leftmargin=1.5em]
  \item Exprimer le couple $C_e$ au niveau des roues d'un essieu en fonction de $F_e$ et de $d_e$. Calculer $C_{e1}$ et $C_{e2}$.
  \item En raisonnant sur les puissances et en tenant compte du rendement $\eta_{red}$, exprimer le couple moteur $C_m$ en fonction de $C_e$, $\eta_{red}$ et $\rho$. Calculer $C_{m1}$ et $C_{m2}$.
  \item En d\'eduire le couple thermique \'equivalent $C_{th}$ sur le cycle de fonctionnement.
  \item Les moteurs \`a courant continu install\'es sont donn\'es pour une puissance nominale $P_{nom} = \SI{96}{\kilo\watt}$ \`a la vitesse nominale $n_{nom} = 1\,035\ \si{\tr\per\minute}$. Conclure sur leur aptitude \`a assurer ce cycle.
\end{enumerate}

}{
\textbf{Corrig\'e.}

\begin{enumerate}[leftmargin=1.5em]
  \item L'effort de traction $F_e$ s'exerce au contact roue/rail, \`a la distance $d_e/2$ de l'axe de l'essieu :
  \[ C_e = F_e \cdot \dfrac{d_e}{2} = F_e \times 0{,}43 \]
  \[ C_{e1} = 14{,}9\times10^{3}\times0{,}43 = 6\,407\ \si{\newton\meter}, \qquad C_{e2} = 9{,}2\times10^{3}\times0{,}43 = 3\,956\ \si{\newton\meter} \]
  \item Le r\'educteur transmet la puissance moteur $P_m = C_m\,\Omega_m$ vers l'essieu avec le rendement $\eta_{red}$ : $C_e\,\Omega_e = \eta_{red}\cdot C_m\,\Omega_m$. Le rapport de r\'eduction \'etant $\rho = \Omega_e/\Omega_m$, il vient
  \[ C_m = \dfrac{C_e \cdot \rho}{\eta_{red}} \]
  \[ C_{m1} = \dfrac{6\,407 \times 0{,}17}{0{,}97} \approx 1\,122{,}9\ \si{\newton\meter}, \qquad C_{m2} = \dfrac{3\,956 \times 0{,}17}{0{,}97} \approx 693{,}3\ \si{\newton\meter} \]
  \item Contributions des deux premi\`eres phases :
  \[ C_{m1}^2\,\Delta t_1 = (1\,122{,}9)^2 \times 42 = 1{,}261\times10^{6}\times42 \approx 52{,}96\times10^{6}\ \si{\newton\squared\meter\squared\second} \]
  \[ C_{m2}^2\,\Delta t_2 = (693{,}3)^2 \times 51 = 0{,}4807\times10^{6}\times51 \approx 24{,}52\times10^{6}\ \si{\newton\squared\meter\squared\second} \]
  Somme sur le cycle complet :
  \[ \sum_i C_{mi}^2\,\Delta t_i = (52{,}96 + 24{,}52 + 558)\times10^{6} \approx 635{,}5\times10^{6}\ \si{\newton\squared\meter\squared\second} \]
  \[ C_{th} = \sqrt{\dfrac{635{,}5\times10^{6}}{2\,048}} = \sqrt{310{,}3\times10^{3}} \approx 557\ \si{\newton\meter} \]
  On note que le terme cumul\'e des seize phases restantes ($558\times10^{6}$) domine largement les deux premi\`eres phases, malgr\'e leurs couples \'elev\'es : ce sont les longues phases de roulage, et non les pointes d'effort, qui d\'eterminent l'\'echauffement.
  \item Couple nominal des moteurs :
  \[ \Omega_{nom} = 1\,035 \times \dfrac{2\pi}{60} \approx \SI{108,4}{\radian\per\second}, \qquad C_{nom} = \dfrac{P_{nom}}{\Omega_{nom}} = \dfrac{96\times10^{3}}{108{,}4} \approx 886\ \si{\newton\meter} \]
  $C_{th} \approx 557\ \si{\newton\meter} < C_{nom} \approx 886\ \si{\newton\meter}$ : les moteurs conviennent du point de vue thermique.
\end{enumerate}

\smallskip
\textbf{Conclusion.} Le couple thermique \'equivalent $C_{th} \approx 557\ \si{\newton\meter}$ repr\'esente $63\,\%$ du couple nominal $C_{nom} \approx 886\ \si{\newton\meter}$ : le crit\`ere de dimensionnement en service intermittent ($C_{th} \leq C_{nom}$, service $S1$ \'equivalent) est satisfait avec une marge de $37\,\%$. C'est bien $C_{th}$, et non le couple de pointe $C_{m1} \approx 1\,123\ \si{\newton\meter}$ -- pourtant sup\'erieur \`a $C_{nom}$ --, qui gouverne l'\'echauffement du bobinage : le d\'epassement de couple sur les $\SI{42}{\second}$ de la phase 1 est thermiquement admissible parce qu'il est bref devant la constante de temps thermique du moteur.
}
\end{tcolorbox}

\CAPETENDFILE
